% !TITLE 把酒问月
% !AUTHOR 李白
% !DYNASTY 唐
% !ORDER 13

\begin{poem}
青天有月来几时？我今停杯一问之。
人攀明月不可得，月行却与人相随。
皎如飞镜临丹阙\textsuperscript{1}，绿烟灭尽\textsuperscript{2}清辉发。
但见宵从海上来，宁知晓向云间没。
白兔捣药\textsuperscript{3}秋复春，嫦娥\textsuperscript{4}孤栖与谁邻？
今人不见古时月，今月曾经照古人。
古人今人若流水，共看明月皆如此。
唯愿当歌对酒时，月光长照金樽里。
\end{poem}

\begin{annotations}
\notetext{1}{丹阙：朱红色的宫阙。月亮像一面飞镜照在宫殿上。}
\notetext{2}{绿烟灭尽：傍晚的烟霭散尽，月色更加清澈。绿烟，黄昏时的青色雾气。}
\notetext{3}{白兔捣药：古人认为月亮中有白兔在捣长生不老之药。年复一年，白兔春秋不停地捣药。}
\notetext{4}{嫦娥：传说嫦娥偷吃了长生药飞到月亮上，独自住在月宫中。“孤栖与谁邻”——她孤独地住在月宫里，和谁做邻居呢？}
\end{annotations}

\begin{appreciation}
这是中国文学史上最著名的问月诗之一。古人写月几千年，多写它的圆缺、清辉，李白不一样，他端起酒杯，直接问月亮：你是什么时候来的？

青天有月来几时？我今停杯一问之。这个起笔随随便便，像聊天时想起个问题。但往深处想，它够吓人——人站在月亮底下几万年，第一次认真问起月亮是什么时候有的，其实就是在问：这世界从什么时候开始的？我在宇宙里算什么？还记得《天问》吗？屈原对着苍天一口气问了一百多个问题。李白的停杯一问，接的正是屈原问天的衣钵。

人攀明月不可得，月行却与人相随。你想攀上月亮，够不着；可你走到哪，月亮跟到哪。可望而不可即——还记得《蒹葭》里“所谓伊人，在水一方”吗？诗人都绕不开这个句式：最想要的，总是隔着一段够不着的距离。

接下来李白数月宫里的事：白兔捣药秋复春，嫦娥孤栖与谁邻？白兔一年年捣着长生药，嫦娥一个人住月宫，连个邻居都没有。长生不老听起来是好事，在李白的想象里，却冷清得可怕——人间虽然短暂，至少热闹。

今人不见古时月，今月曾经照古人。这是全诗最出名的两句。几乎同时，张若虚在《春江花月夜》里问“江畔何人初见月？江月何年初照人”，问的是同一个问题：月亮看了多少人，人却只看了一小会儿月亮。古人今人若流水，共看明月皆如此——古人和今人像流水一波波过去，月亮还是同一轮。这么一比，人的一生实在短得可怜。

但李白没有停在悲伤里。最后他回到眼前：唯愿当歌对酒时，月光长照金樽里。问了一圈宇宙的大问题，答案落在自己的酒杯里——人生有限，月亮恒久，那就把今晚这杯酒，在月亮底下喝好。后来苏轼写“明月几时有？把酒问青天”，问的就是李白这一句。这本书往后读，月亮还会出现很多次，从《静夜思》照到这一首，再照到苏轼的水调歌头，一直亮着。

月亮不会回答任何问题。但下次你抬头看它，可以想一想：这轮月亮照过李白，照过苏轼，也照着咱们家的阳台。人换了一茬又一茬，月亮没变过——这么一想，眼下的烦恼也就没那么大了。看完早点睡，月亮明天还在。
\end{appreciation}
